% WRITER: Mahek. Three paragraphs:
% (1) Agentic benchmarks: SWE-bench [jimenez2024swebench], Terminal-Bench [merrill2026terminalbench, tb4],
%     SWE-Bench Pro [deng2025sweprobench]; harness effects: SWE-agent [yang2024sweagent], OpenHands [wang2025openhands].
% (2) Routing and cascades: RouteLLM [ong2024routellm], FrugalGPT [chen2023frugalgpt], AutoMix (add bib entry),
%     Arch-Router [tran2025archrouter], NVIDIA classifier [nvidia2026classifier], Laya [laya2026].
% (3) Learning to defer / selective prediction; multi-agent coordination: MetaGPT [hong2024metagpt].
% Position our work: routers decide from the prompt; our cascade decides after the cheap model's attempt.
\section{Related Work}
\label{sec:related}

\paragraph{Agentic benchmarks and harness effects.}
SWE-bench~\cite{jimenez2024swebench} evaluates whether a model can resolve real GitHub issues: it collects 2,294
issues from 12 Python repositories and checks each proposed patch against the repository's tests. SWE-Bench
Pro~\cite{deng2025sweprobench} extends this design to 1,865 longer problems from 41 actively maintained
repositories, many of which need changes across several files, and keeps part of its data private to limit
contamination. Terminal-Bench~\cite{merrill2026terminalbench} moves the setting to the command line: each task
has its own environment, a human-written reference solution and tests that verify the result; we use its
version~4.0~\cite{tb4}. Work on agent design has shown that the software around the model matters.
SWE-agent~\cite{yang2024sweagent} found that the interface through which a model reads and edits code changes
how many issues the same model resolves, and OpenHands~\cite{wang2025openhands} provides an open platform in which
agents act through a sandboxed shell, code editor and browser. These studies mainly design and improve one
harness. Our Phase~1 takes the opposite view: it holds the model fixed and compares four existing
harnesses, for three models, on the same tasks.

\paragraph{Routing and cascades.}
Model routing aims to send each query to the least expensive model that can handle it. Most routers decide
before any work is done, from the query text alone. RouteLLM~\cite{ong2024routellm} learns from human preference
data whether a query needs a strong or a weak model. Arch-Router~\cite{tran2025archrouter}, a 1.5B-parameter
model, maps each query to user-defined domains and action types and routes by those labels. NVIDIA's prompt
classifier~\cite{nvidia2026classifier} assigns a task type and scores several dimensions of complexity, such as
reasoning and domain knowledge, and Laya~\cite{laya2026} returns typed decisions that we used to score the
capability each step needed in our first design. Cascades instead decide after a model has answered.
FrugalGPT~\cite{chen2023frugalgpt} queries a sequence of models and stops once a learned scoring function
accepts an answer, and AutoMix~\cite{aggarwal2024automix} has the smaller model check its own answer and uses a
router that accounts for the noise in that check to decide whether to call a larger model. Both were designed
for queries with a single text answer. In our setting the object being judged is the state of a working
environment after many commands, and we found that routing from task text did not separate tasks by required
capability (Section~\ref{sec:p2method}). Our cascade therefore keeps the decision until the inexpensive model
has produced something the strong model can inspect.

\paragraph{Learning to defer and multi-agent coordination.}
Learning to defer trains a model to either make a prediction or pass the case to an expert, so that the
combined system performs better~\cite{madras2018predict,mozannar2020consistent}. It is
closely related to selective prediction, in which a model abstains when it is not confident
enough~\cite{geifman2017selective}. Our cascade applies the same idea to agents: the inexpensive model always
acts first, and a stronger model reviews the result and decides whether to take over. In the standard formulation the
decision to defer is made from the input; here it is made from the outcome of an attempt, and the expert
continues from that attempt rather than starting again. Multi-agent frameworks such as MetaGPT~\cite{hong2024metagpt} divide a software task among
agents with fixed roles that pass structured documents to each other. Our first design followed a similar plan,
with a planner and per-step workers. On Terminal-Bench, however, steps shared hidden state such as files and
values discovered along the way, so splitting a task added hand-offs in which information was lost.

